\documentclass[a4paper]{article}     
\usepackage{amsfonts}
\usepackage{amsmath}
\usepackage{amssymb}
\usepackage{graphicx}
\DeclareMathOperator{\cis}{cis}

\begin{document}
	
	\noindent \large Auteur: Victor Pena \\
	\noindent \large Studierichting: Informatica\\
	\noindent \large Oefeningenreeks 1D nr. 21.4\\
	
	\medskip
	
	\normalsize 
	
	$P(z)=z^3 +3z^2 + 12z + 10$ \\
	
	
	$P(-1)=-1 + 3 -12 +10 = 0$ \\

	$\Rightarrow$ \begin{tabular}{l|lll|l}
		 & $1$   & $3$  & $12$ &$10$ \\
	$-1$ &       & $-1$ & $-2$ &$-10$ \\  \cline{1-5}
		 & $1$   & $2 $ & $10$ &$0$
	\end{tabular} $\Rightarrow P(z) = (z+1)(z^2 + 2z + 10)$ \\
	

	De wortels van $z^2 + 2z + 10$ zijn gegeven door: \\
	
	$z = \dfrac{-2 \pm i\sqrt{40-4}}{2} = \dfrac{-2 \pm 6i}{2} = -1 \pm 3i$\\
	
	
	Dus, $P(z) = (z+1)(z+1-3i)(z+1+3i)$, met wortels $-1,-1+3i$ en $-1-3i$ \\
	
	
	
	
	
	
	
	
	
	
	
	
	\begin{thebibliography}{999}
		%\bibitem{a} Referentie
	\end{thebibliography}
\end{document} 
